% GENERATED FILE. Edit canonical lessons, then run npm run generate:books.
% canonical-source-hash: fnv1a64:6e02038a0980f628
% canonical-lessons: RU-C522-skulptor, RU-C522-skulptura, RU-C522-mramor, RU-C522-statuya, RU-C522-vdokhnovlyat

\chapter{Marble and Statues}
\label{ch:ru-marble-and-statues}
\hlchaptermodality{522}

\begin{chapteropening}
\textbf{By the end of this chapter:} \emph{I can talk about sculpture and what inspires artists.}

Complete the last lesson of chapter 522: I can talk about sculpture and what inspires artists.
\end{chapteropening}

\section[\texorpdfstring{\ru{скульптор} (skúlptor)}{скульптор (skúlptor)}]{\ru{скульптор} (skúlptor) — a sculptor}

\label{lesson:RU-C522-skulptor}



\begin{quote}
Before the new one: say the Russian for a landscape, then the Russian for to pose.
\end{quote}

\subsection*{You'll want to know: \ru{скульптор}}
\textbf{\ru{скульптор}} (\emph{skúlptor}) — \textquotedblleft{}a sculptor\textquotedblright{}.

\begin{quote}
\ru{Скульптор} \ru{работал} \ru{над} \ru{памятником} \ru{целый} \ru{год}.
\end{quote}

\begin{quote}
\emph{Skúlptor rabótal nad pámyatnikom tsélyy god.}
\end{quote}

\begin{quote}
The sculptor worked on the monument for a whole year.
\end{quote}

Masculine.

\begin{cousinweb}
\emph{\ru{скульптор}} is a learned borrowing from Latin \emph{sculptor}, built on \emph{sculpere} 'to carve' (Lewis \& Short); EtymDB has no entry.
\end{cousinweb}

\begin{grammarlens}[title={What you've built}]
One more word for telling, explaining and describing at B1.
\end{grammarlens}

\subsection*{Your turn}
\begin{itemize}
\raggedright
  \item \emph{Say it:} \emph{skúlptor}
  \item \emph{Say it:} \emph{skúlptor}, once more
  \item \emph{Say it:} a sentence of your own with \emph{skúlptor}
  \item \emph{Recall:} say the Russian for a landscape, then the Russian for to pose, then say \emph{skúlptor} again
\end{itemize}

\begin{tcolorbox}[breakable,colback=teal!4,colframe=teal!35!black,title={Before you move on}]
What does \ru{скульптор} mean? (\textquotedblleft{}A sculptor\textquotedblright{}.) Say it once more.
\end{tcolorbox}
\section[\texorpdfstring{\ru{скульптура} (skulptúra)}{скульптура (skulptúra)}]{\ru{скульптура} (skulptúra) — a sculpture}

\label{lesson:RU-C522-skulptura}



\begin{quote}
Before the new one: say the Russian for to pose, then the Russian for a sculptor.
\end{quote}

\subsection*{You'll want to know: \ru{скульптура}}
\textbf{\ru{скульптура}} (\emph{skulptúra}) — \textquotedblleft{}a sculpture\textquotedblright{}.

\begin{quote}
\ru{В} \ru{музее} \ru{много} \ru{древних} \ru{скульптур}.
\end{quote}

\begin{quote}
\emph{V muzéye mnógo drévnikh skulptúr.}
\end{quote}

\begin{quote}
There are many ancient sculptures in the museum.
\end{quote}

Feminine.

\begin{cousinweb}
\emph{\ru{скульптура}} is a learned borrowing from Latin \emph{sculptūra} 'carving', from \emph{sculpere} 'to carve' (Lewis \& Short); EtymDB lists the Russian word without a source.
\end{cousinweb}

\begin{grammarlens}[title={What you've built}]
One more word for telling, explaining and describing at B1.
\end{grammarlens}

\subsection*{Your turn}
\begin{itemize}
\raggedright
  \item \emph{Say it:} \emph{skulptúra}
  \item \emph{Say it:} \emph{skulptúra}, once more
  \item \emph{Say it:} a sentence of your own with \emph{skulptúra}
  \item \emph{Recall:} say the Russian for to pose, then the Russian for a sculptor, then say \emph{skulptúra} again
\end{itemize}

\begin{tcolorbox}[breakable,colback=teal!4,colframe=teal!35!black,title={Before you move on}]
What does \ru{скульптура} mean? (\textquotedblleft{}A sculpture\textquotedblright{}.) Say it once more.
\end{tcolorbox}
\section[\texorpdfstring{\ru{мрамор} (mrámor)}{мрамор (mrámor)}]{\ru{мрамор} (mrámor) — marble}

\label{lesson:RU-C522-mramor}



\begin{quote}
Before the new one: say the Russian for a sculptor, then the Russian for a sculpture.
\end{quote}

\subsection*{You'll want to know: \ru{мрамор}}
\textbf{\ru{мрамор}} (\emph{mrámor}) — \textquotedblleft{}marble\textquotedblright{}.

\begin{quote}
\ru{Лестница} \ru{во} \ru{дворце} \ru{сделана} \ru{из} \ru{мрамора}.
\end{quote}

\begin{quote}
\emph{Léstnitsa vo dvortsé sdélana iz mrámora.}
\end{quote}

\begin{quote}
The staircase in the palace is made of marble.
\end{quote}

Masculine.

\begin{cousinweb}
\emph{\ru{мрамор}} was borrowed from Church Slavonic \emph{mramor\ru{ъ}}, the \emph{-\ru{ра}-} shape of Proto-Slavic \emph{mormor\ru{ъ}}, which came from Latin \emph{marmor} and Greek \emph{mármaros} (EtymDB).
\end{cousinweb}

\begin{grammarlens}[title={What you've built}]
One more word for telling, explaining and describing at B1.
\end{grammarlens}

\subsection*{Your turn}
\begin{itemize}
\raggedright
  \item \emph{Say it:} \emph{mrámor}
  \item \emph{Say it:} \emph{mrámor}, once more
  \item \emph{Say it:} a sentence of your own with \emph{mrámor}
  \item \emph{Recall:} say the Russian for a sculptor, then the Russian for a sculpture, then say \emph{mrámor} again
\end{itemize}

\begin{tcolorbox}[breakable,colback=teal!4,colframe=teal!35!black,title={Before you move on}]
What does \ru{мрамор} mean? (\textquotedblleft{}Marble\textquotedblright{}.) Say it once more.
\end{tcolorbox}
\section[\texorpdfstring{\ru{статуя} (státuya)}{статуя (státuya)}]{\ru{статуя} (státuya) — a statue}

\label{lesson:RU-C522-statuya}



\begin{quote}
Before the new one: say the Russian for a sculpture, then the Russian for marble.
\end{quote}

\subsection*{You'll want to know: \ru{статуя}}
\textbf{\ru{статуя}} (\emph{státuya}) — \textquotedblleft{}a statue\textquotedblright{}.

\begin{quote}
\ru{В} \ru{парке} \ru{стоит} \ru{статуя} \ru{из} \ru{бронзы}.
\end{quote}

\begin{quote}
\emph{V párke stoít státuya iz brónzy.}
\end{quote}

\begin{quote}
There is a bronze statue in the park.
\end{quote}

Feminine.

\begin{cousinweb}
\emph{\ru{статуя}} is a learned borrowing from Latin \emph{statua} 'an image, a statue', which Lewis \& Short derive from \emph{statuere} 'to set up'; EtymDB has no entry.
\end{cousinweb}

\begin{grammarlens}[title={What you've built}]
One more word for telling, explaining and describing at B1.
\end{grammarlens}

\subsection*{Your turn}
\begin{itemize}
\raggedright
  \item \emph{Say it:} \emph{státuya}
  \item \emph{Say it:} \emph{státuya}, once more
  \item \emph{Say it:} a sentence of your own with \emph{státuya}
  \item \emph{Recall:} say the Russian for a sculpture, then the Russian for marble, then say \emph{státuya} again
\end{itemize}

\begin{tcolorbox}[breakable,colback=teal!4,colframe=teal!35!black,title={Before you move on}]
What does \ru{статуя} mean? (\textquotedblleft{}A statue\textquotedblright{}.) Say it once more.
\end{tcolorbox}
\section[\texorpdfstring{\ru{вдохновлять} (vdokhnovlyát)}{вдохновлять (vdokhnovlyát)}]{\ru{вдохновлять} (vdokhnovlyát) — to inspire}

\label{lesson:RU-C522-vdokhnovlyat}



\begin{quote}
Before the new one: say the Russian for marble, then the Russian for a statue.
\end{quote}

\subsection*{You'll want to know: \ru{вдохновлять}}
\textbf{\ru{вдохновлять}} (\emph{vdokhnovlyát}) — \textquotedblleft{}to inspire\textquotedblright{}.

\begin{quote}
\ru{Море} \ru{вдохновляет} \ru{художников}.
\end{quote}

\begin{quote}
\emph{Móre vdokhnovlyáyet khudózhnikov.}
\end{quote}

\begin{quote}
The sea inspires artists.
\end{quote}

Imperfective; perfective \ru{вдохновить} (\ru{вдохновлю}). The noun is \emph{\ru{вдохновение}} — inspiration.

\begin{cousinweb}
\emph{\ru{вдохновлять}} is built on \emph{\ru{вдохнуть}} 'to breathe in', from the root of \emph{\ru{дух}} 'breath, spirit', which EtymDB traces to Proto-Slavic \emph{dux\ru{ъ}}.
\end{cousinweb}

\begin{grammarlens}[title={What you've built}]
One more word for telling, explaining and describing at B1.
\end{grammarlens}

\subsection*{Your turn}
\begin{itemize}
\raggedright
  \item \emph{Say it:} \emph{vdokhnovlyát}
  \item \emph{Say it:} \emph{vdokhnovlyát}, once more
  \item \emph{Say it:} a sentence of your own with \emph{vdokhnovlyát}
  \item \emph{Recall:} say the Russian for marble, then the Russian for a statue, then say \emph{vdokhnovlyát} again
\end{itemize}

\begin{tcolorbox}[breakable,colback=teal!4,colframe=teal!35!black,title={Before you move on}]
What does \ru{вдохновлять} mean? (\textquotedblleft{}To inspire\textquotedblright{}.) Say it once more.
\end{tcolorbox}
